% 不等式约束可行域：等高线 + 多个约束 + 阴影可行域 + 边界最优解
\documentclass[tikz,border=8pt]{standalone}
\usepackage{ctex}
\usepackage{amsmath}
\usepackage{pgfplots}
\pgfplotsset{compat=1.18}
\usetikzlibrary{calc,arrows.meta,patterns,intersections}

\begin{document}
\begin{tikzpicture}[scale=1.0]

%% 坐标轴
\draw[->, gray!60, thin] (-0.5,0) -- (5.5,0) node[right, font=\small] {$x_1$};
\draw[->, gray!60, thin] (0,-0.5) -- (0,4.8) node[above, font=\small] {$x_2$};

%% 可行域阴影（交集：x1>=0, x2>=0, x1+x2<=4, x1<=3）
%% 顶点：(0,0),(3,0),(3,1),(0,4)（从 x1+x2=4 和 x1=3 交点 (3,1)）
\fill[blue!15] (0,0) -- (3,0) -- (3,1) -- (0,4) -- cycle;

%% 约束边界线
%% g1: x1 + x2 = 4（直线，可行域在下方）
\draw[blue!70!black, thick, name path=g1line]
  (-0.2, 4.2) -- (4.2, -0.2);
\node[blue!70!black, font=\small, rotate=-45] at (3.9, 0.5) {$g_1: x_1{+}x_2{\leq}4$};

%% g2: x1 = 3（竖直线，可行域在左方）
\draw[red!70!black, thick]
  (3, -0.3) -- (3, 4.5);
\node[red!70!black, font=\small, rotate=90] at (3.25, 3.5) {$g_2: x_1{\leq}3$};

%% 非负约束（坐标轴本身）
\node[gray!70, font=\scriptsize] at (1.0, -0.3) {$x_1 \geq 0$};
\node[gray!70, font=\scriptsize] at (-0.45, 1.5) [rotate=90] {$x_2 \geq 0$};

%% 可行域标注
\node[blue!60!black, font=\small\bfseries] at (1.2, 1.5) {可行域 $\mathcal{F}$};

%% 等高线族：f(x1,x2) = (x1-2)^2 + (x2-3)^2（目标函数，圆心在(2,3)）
\foreach \r in {0.5, 1.0, 1.5, 2.2, 3.0} {
  \draw[gray!45, thin] (2,3) circle (\r);
}
\node[gray!60, font=\scriptsize] at (2.55, 3.08) {$f{=}c_1$};
\node[gray!60, font=\scriptsize] at (3.1, 3.08) {$f{=}c_2$};
\node[gray!60, font=\scriptsize] at (3.6, 3.08) {$f{=}c_3$};

%% 无约束最优（圆心）(2,3)：不在可行域内（x1+x2=5>4）
\fill[gray!60] (2,3) circle (2pt);
\node[gray!60, font=\scriptsize, above right] at (2,3) {无约束最优（不可行）};

%% 带约束最优点：在边界 x1+x2=4 上
%% 最优：从 (2,3) 垂直投影到直线 x1+x2=4 => 投影点 = (2+(4-5)/2, 3+(4-5)/2) = (1.5, 2.5)
\coordinate (Pstar) at (1.5, 2.5);
\fill[red!80!black] (Pstar) circle (3pt);
\node[red!80!black, font=\small\bfseries, above left] at (Pstar)
  {$\mathbf{x}^*$（活跃约束 $g_1$）};

%% 活跃约束高亮
\draw[red!80!black, very thick, dashed]
  (0, 4) -- (3, 1);
\node[red!80!black, font=\scriptsize] at (0.6, 3.8) {活跃约束边界};

%% 在最优点处画 nabla f（指向可行域外）
\draw[->, very thick, red!80!black, >=Stealth]
  (Pstar) -- ++(0.5, 0.5) node[right, font=\small] {$\nabla f$};

%% 在最优点处画 nabla g1（与 nabla f 共线，反向）
\draw[->, very thick, blue!70!black, >=Stealth]
  (Pstar) -- ++(-0.35, -0.35) node[below left, font=\small, blue!70!black] {$\nabla g_1$};

%% 标注 KKT 条件在最优点
\node[draw=gray!50, fill=white, rounded corners, font=\scriptsize, inner sep=3pt,
      align=left] at (4.2, 3.8)
  {KKT 最优点：\\$\nabla f + \mu_1\nabla g_1 = \mathbf{0}$\\$\mu_1>0$（$g_1$ 活跃）\\$\mu_2=0$（$g_2$ 非活跃）};

%% 顶点标注
\fill (0,0) circle (2pt) node[below left, font=\scriptsize] {$(0,0)$};
\fill (3,0) circle (2pt) node[below, font=\scriptsize] {$(3,0)$};
\fill (3,1) circle (2pt) node[right, font=\scriptsize] {$(3,1)$};
\fill (0,4) circle (2pt) node[left, font=\scriptsize] {$(0,4)$};

%% 标题
\node[font=\small\bfseries, align=center] at (2.5, -0.8)
  {不等式约束可行域（阴影）、等高线与边界最优解};

\end{tikzpicture}
\end{document}
